% =====================================================================
% Informe maestro - Parte 1 (Kamila): mision, stakeholders, objetivos y alcance
% Corrige y completa lo escrito en la semana 2 (secciones 2.1, 2.2 y 3.1 de aquel informe).
% Es un FRAGMENTO: no tiene \documentclass ni \begin{document}.
% Para unirlo al informe: % =====================================================================
% Informe maestro - Parte 1 (Kamila): mision, stakeholders, objetivos y alcance
% Corrige y completa lo escrito en la semana 2 (secciones 2.1, 2.2 y 3.1 de aquel informe).
% Es un FRAGMENTO: no tiene \documentclass ni \begin{document}.
% Para unirlo al informe: % =====================================================================
% Informe maestro - Parte 1 (Kamila): mision, stakeholders, objetivos y alcance
% Corrige y completa lo escrito en la semana 2 (secciones 2.1, 2.2 y 3.1 de aquel informe).
% Es un FRAGMENTO: no tiene \documentclass ni \begin{document}.
% Para unirlo al informe: % =====================================================================
% Informe maestro - Parte 1 (Kamila): mision, stakeholders, objetivos y alcance
% Corrige y completa lo escrito en la semana 2 (secciones 2.1, 2.2 y 3.1 de aquel informe).
% Es un FRAGMENTO: no tiene \documentclass ni \begin{document}.
% Para unirlo al informe: \input{informe_parte1}
% Paquetes que usa: booktabs, longtable, array, enumitem (ya van en el preambulo habitual).
% =====================================================================

\section{Misión}

El \textit{Repositorio de Tutorías} existe para guardar en un solo lugar, y clasificados
por periodo académico y por categoría documental, los documentos de las tutorías: el plan
de acción inicial, el informe de corte y el informe final.

Hoy esa información está repartida en correos electrónicos, archivos sueltos y
dispositivos personales. Cuando a un coordinador le piden un histórico de cinco años,
tiene que revisar cinco años de correos, y al no existir un lugar centralizado se pierde
información y cuesta encontrarla. Con la plataforma basta con cambiar de periodo y
consultar lo que se subió en ese semestre, en cualquier momento y desde cualquier lugar.

La plataforma no está pensada solo para la carrera de Software. Las facultades y las
carreras se registran como datos del sistema, de modo que cualquier facultad y cualquier
carrera de la institución pueda usarla con su propia información.

\section{Stakeholders}

La Tabla~\ref{tab:stakeholders} indica quiénes tienen interés en la plataforma o se ven
afectados por ella, y cuál es su relación con el sistema. Los tutores son quienes suben
los documentos; los demás roles organizan, consultan o evalúan.

\begin{longtable}{>{\raggedright\arraybackslash}p{3.7cm}>{\raggedright\arraybackslash}p{11.3cm}}
\caption{Stakeholders del sistema y su relación con él}\label{tab:stakeholders}\\
\toprule
\textbf{Stakeholder} & \textbf{Relación con el sistema} \\
\midrule
\endfirsthead
\toprule
\textbf{Stakeholder} & \textbf{Relación con el sistema} \\
\midrule
\endhead
Administrador general &
Configura la plataforma: registra facultades, carreras y usuarios, crea los periodos
y define qué categorías documentales se piden en cada semestre (por ejemplo, plan de
acción inicial, informe de corte e informe final). \\
\addlinespace
Coordinador de carrera &
Da seguimiento a los documentos de su carrera. Consulta lo que entregaron los tutores,
filtrando por periodo y categoría. \\
\addlinespace
Coordinador de área &
Revisa los documentos de las carreras de su área. \\
\addlinespace
Tutores &
Son quienes suben los documentos. Eligen el periodo y la categoría, escriben el título
y adjuntan el archivo, que queda guardado en PostgreSQL. \\
\addlinespace
Docente de la materia (evaluador) &
Orienta y evalúa el proyecto. Revisa que la aplicación funcione, que la integración con
PostgreSQL sea correcta y que las decisiones estén documentadas. \\
\addlinespace
Institución (UTEQ) &
Es la beneficiaria final: conserva el histórico de tutorías ordenado y accesible,
sin depender de correos personales. \\
\bottomrule
\end{longtable}

\section{Objetivos}

\subsection{Objetivo general}

Desarrollar una aplicación web conectada a PostgreSQL que guarde y clasifique por
periodo académico y categoría documental los documentos de tutorías, para que pueda
ser usada por cualquier facultad y carrera y reemplace la búsqueda de documentos en
años de correos.

\subsection{Objetivos específicos}

\begin{enumerate}[label=\textbf{OE\arabic*.}, leftmargin=1.6cm]
  \item Gestionar periodos y categorías documentales, de modo que al crear un periodo
        nuevo se puedan reutilizar las categorías de periodos anteriores y agregar
        nuevas o quitar las que ya no se necesiten (módulo de periodos y categorías,
        RF-01).
  \item Registrar facultades, carreras y usuarios, asignando a cada usuario uno de los
        cuatro roles (administrador general, coordinador de carrera, coordinador de área
        o tutor), con validaciones de integridad como nombre único, correo válido y
        único, y carrera obligatoria para todos los roles excepto el administrador
        general (módulo de facultades, carreras y usuarios, RF-02).
  \item Permitir que los tutores suban documentos y que se consulten con filtros por
        periodo, categoría y carrera, con paginación, validando el tipo de archivo, el
        tamaño y que la categoría esté habilitada para ese periodo (módulo de
        documentos, RF-03).
  \item Diseñar la base de datos para que soporte millones de registros, usando índices
        en las llaves foráneas y en los campos de filtro, paginación en las listas y un
        script de datos de prueba masivos, para que las pantallas no se vuelvan lentas
        al crecer la información.
  \item Proteger la integridad de los datos con llaves foráneas \texttt{ON DELETE
        RESTRICT}, consultas parametrizadas y mensajes claros de éxito o de error en
        cada operación.
  \item Comprobar cada requisito con pruebas funcionales que registren los datos
        usados, el resultado esperado, el resultado obtenido y su evidencia, y que
        verifiquen en PostgreSQL lo que se hizo desde la aplicación.
\end{enumerate}

\section{Alcance}

\subsection{Qué incluye este semestre}

La primera versión de la aplicación es una aplicación web en C\# con ASP.NET Core MVC
y PostgreSQL como gestor de base de datos, donde se almacena toda la información. Incluye
tres módulos de prioridad alta (Tabla~\ref{tab:modulos}):

\begin{longtable}{>{\raggedright\arraybackslash}p{3.7cm}>{\raggedright\arraybackslash}p{11.3cm}}
\caption{Módulos incluidos en la primera versión}\label{tab:modulos}\\
\toprule
\textbf{Módulo} & \textbf{Funcionalidades incluidas} \\
\midrule
\endfirsthead
\toprule
\textbf{Módulo} & \textbf{Funcionalidades incluidas} \\
\midrule
\endhead
Periodos y categorías documentales (RF-01) &
Crear, listar, editar, cerrar y eliminar periodos; crear, editar y eliminar categorías
documentales; reutilizar categorías de periodos anteriores al crear un periodo nuevo.
Validaciones: nombre de periodo único, fecha de fin posterior a la de inicio, categoría
con nombre y no eliminar una categoría que ya tiene documentos. \\
\addlinespace
Facultades, carreras y usuarios (RF-02) &
Crear, listar, editar y eliminar facultades, carreras (cada una pertenece a una facultad)
y usuarios (con rol y carrera). Validaciones: nombre único, correo con formato válido y
único, carrera con facultad, carrera obligatoria salvo para el administrador general, y
no eliminar una facultad con carreras ni un usuario con documentos. \\
\addlinespace
Documentos (RF-03) &
Subir documentos en PDF o DOCX, consultarlos con filtros por periodo, categoría y carrera
con paginación y orden por fecha, editar sus datos, eliminarlos y descargarlos. \\
\bottomrule
\end{longtable}

Además, el alcance incluye la conexión de la aplicación con PostgreSQL, los scripts SQL
por requisito (\texttt{00\_base.sql}, \texttt{rf01.sql}, \texttt{rf02.sql} y
\texttt{rf03.sql}), las pruebas funcionales con evidencia y un script de datos de prueba
masivos para comprobar el rendimiento del diseño.

\subsection{Qué queda fuera}

\begin{itemize}[leftmargin=1.2cm]
  \item Inicio de sesión y control de acceso, que se contemplan para una versión posterior.
        Por ahora los roles se registran en el sistema, pero todavía no se restringen las
        pantallas según el rol de quien entra.
  \item Administración avanzada de roles y privilegios de PostgreSQL, que corresponde a
        contenidos posteriores de la materia.
  \item Carga real de millones de documentos. Solo se generan datos ficticios para
        probar que el diseño responde bien.
  \item Videollamadas, mensajería entre usuarios y agendamiento de tutorías, porque el
        problema que se resuelve es la dispersión y la pérdida de los documentos.
  \item Reportes y estadísticas de uso de la plataforma.
  \item Avisos o recordatorios por correo a los tutores.
  \item Importar automáticamente los documentos históricos que hoy están en correos.
\end{itemize}

\subsection{Cambios respecto a la propuesta inicial}

Después de la propuesta de la semana 2 el docente precisó cómo debe funcionar la
plataforma. Por eso se ajustó la propuesta inicial. La Tabla~\ref{tab:cambios} resume
los cambios y su motivo.

\begingroup\small
\begin{longtable}{>{\raggedright\arraybackslash}p{2.3cm}>{\raggedright\arraybackslash}p{3.6cm}>{\raggedright\arraybackslash}p{4.3cm}>{\raggedright\arraybackslash}p{3.9cm}}
\caption{Cambios respecto a la propuesta de la semana 2}\label{tab:cambios}\\
\toprule
\textbf{Elemento} & \textbf{Propuesta inicial} & \textbf{Versión actual} & \textbf{Motivo} \\
\midrule
\endfirsthead
\toprule
\textbf{Elemento} & \textbf{Propuesta inicial} & \textbf{Versión actual} & \textbf{Motivo} \\
\midrule
\endhead
Problema &
Materiales de tutorías para estudiantes y docentes. &
Documentos de tutorías (plan de acción inicial, informe de corte e informe final)
que hoy se buscan en años de correos. &
El docente explicó que el objetivo es evitar revisar cinco años de correos. \\
\addlinespace
Actores &
Administrador, docente y estudiante (tres actores). &
Administrador general, coordinador de carrera, coordinador de área y tutor (cuatro
roles). El estudiante ya no es un actor. &
El docente indicó estos roles; los tutores son quienes suben los documentos. \\
\addlinespace
Cobertura &
No se precisaba. &
Cualquier facultad y carrera, con facultades y carreras registradas como datos. &
El docente pidió que no sirva solo a la carrera de Software. \\
\addlinespace
Organización de los documentos &
Por asignatura, tema, tipo y fecha. &
Por periodo académico y categoría documental, reutilizable entre semestres. &
El administrador define cada semestre qué documentos se necesitan y las categorías
se reutilizan. \\
\addlinespace
Seguridad &
Inicio de sesión, control de acceso, cierre de sesión y contraseñas cifradas. &
Se deja para una versión posterior. Por ahora solo se registra el rol de cada usuario. &
Decisión del equipo para priorizar los tres módulos; la administración avanzada de
roles corresponde a contenidos posteriores. \\
\addlinespace
Rendimiento &
Búsquedas ágiles con al menos 1\,000 recursos. &
Diseño pensado para millones de registros, con índices, paginación y datos de prueba
masivos. &
El docente indicó que las bases de datos de la materia manejan millones de registros. \\
\bottomrule
\end{longtable}
\endgroup

% Paquetes que usa: booktabs, longtable, array, enumitem (ya van en el preambulo habitual).
% =====================================================================

\section{Misión}

El \textit{Repositorio de Tutorías} existe para guardar en un solo lugar, y clasificados
por periodo académico y por categoría documental, los documentos de las tutorías: el plan
de acción inicial, el informe de corte y el informe final.

Hoy esa información está repartida en correos electrónicos, archivos sueltos y
dispositivos personales. Cuando a un coordinador le piden un histórico de cinco años,
tiene que revisar cinco años de correos, y al no existir un lugar centralizado se pierde
información y cuesta encontrarla. Con la plataforma basta con cambiar de periodo y
consultar lo que se subió en ese semestre, en cualquier momento y desde cualquier lugar.

La plataforma no está pensada solo para la carrera de Software. Las facultades y las
carreras se registran como datos del sistema, de modo que cualquier facultad y cualquier
carrera de la institución pueda usarla con su propia información.

\section{Stakeholders}

La Tabla~\ref{tab:stakeholders} indica quiénes tienen interés en la plataforma o se ven
afectados por ella, y cuál es su relación con el sistema. Los tutores son quienes suben
los documentos; los demás roles organizan, consultan o evalúan.

\begin{longtable}{>{\raggedright\arraybackslash}p{3.7cm}>{\raggedright\arraybackslash}p{11.3cm}}
\caption{Stakeholders del sistema y su relación con él}\label{tab:stakeholders}\\
\toprule
\textbf{Stakeholder} & \textbf{Relación con el sistema} \\
\midrule
\endfirsthead
\toprule
\textbf{Stakeholder} & \textbf{Relación con el sistema} \\
\midrule
\endhead
Administrador general &
Configura la plataforma: registra facultades, carreras y usuarios, crea los periodos
y define qué categorías documentales se piden en cada semestre (por ejemplo, plan de
acción inicial, informe de corte e informe final). \\
\addlinespace
Coordinador de carrera &
Da seguimiento a los documentos de su carrera. Consulta lo que entregaron los tutores,
filtrando por periodo y categoría. \\
\addlinespace
Coordinador de área &
Revisa los documentos de las carreras de su área. \\
\addlinespace
Tutores &
Son quienes suben los documentos. Eligen el periodo y la categoría, escriben el título
y adjuntan el archivo, que queda guardado en PostgreSQL. \\
\addlinespace
Docente de la materia (evaluador) &
Orienta y evalúa el proyecto. Revisa que la aplicación funcione, que la integración con
PostgreSQL sea correcta y que las decisiones estén documentadas. \\
\addlinespace
Institución (UTEQ) &
Es la beneficiaria final: conserva el histórico de tutorías ordenado y accesible,
sin depender de correos personales. \\
\bottomrule
\end{longtable}

\section{Objetivos}

\subsection{Objetivo general}

Desarrollar una aplicación web conectada a PostgreSQL que guarde y clasifique por
periodo académico y categoría documental los documentos de tutorías, para que pueda
ser usada por cualquier facultad y carrera y reemplace la búsqueda de documentos en
años de correos.

\subsection{Objetivos específicos}

\begin{enumerate}[label=\textbf{OE\arabic*.}, leftmargin=1.6cm]
  \item Gestionar periodos y categorías documentales, de modo que al crear un periodo
        nuevo se puedan reutilizar las categorías de periodos anteriores y agregar
        nuevas o quitar las que ya no se necesiten (módulo de periodos y categorías,
        RF-01).
  \item Registrar facultades, carreras y usuarios, asignando a cada usuario uno de los
        cuatro roles (administrador general, coordinador de carrera, coordinador de área
        o tutor), con validaciones de integridad como nombre único, correo válido y
        único, y carrera obligatoria para todos los roles excepto el administrador
        general (módulo de facultades, carreras y usuarios, RF-02).
  \item Permitir que los tutores suban documentos y que se consulten con filtros por
        periodo, categoría y carrera, con paginación, validando el tipo de archivo, el
        tamaño y que la categoría esté habilitada para ese periodo (módulo de
        documentos, RF-03).
  \item Diseñar la base de datos para que soporte millones de registros, usando índices
        en las llaves foráneas y en los campos de filtro, paginación en las listas y un
        script de datos de prueba masivos, para que las pantallas no se vuelvan lentas
        al crecer la información.
  \item Proteger la integridad de los datos con llaves foráneas \texttt{ON DELETE
        RESTRICT}, consultas parametrizadas y mensajes claros de éxito o de error en
        cada operación.
  \item Comprobar cada requisito con pruebas funcionales que registren los datos
        usados, el resultado esperado, el resultado obtenido y su evidencia, y que
        verifiquen en PostgreSQL lo que se hizo desde la aplicación.
\end{enumerate}

\section{Alcance}

\subsection{Qué incluye este semestre}

La primera versión de la aplicación es una aplicación web en C\# con ASP.NET Core MVC
y PostgreSQL como gestor de base de datos, donde se almacena toda la información. Incluye
tres módulos de prioridad alta (Tabla~\ref{tab:modulos}):

\begin{longtable}{>{\raggedright\arraybackslash}p{3.7cm}>{\raggedright\arraybackslash}p{11.3cm}}
\caption{Módulos incluidos en la primera versión}\label{tab:modulos}\\
\toprule
\textbf{Módulo} & \textbf{Funcionalidades incluidas} \\
\midrule
\endfirsthead
\toprule
\textbf{Módulo} & \textbf{Funcionalidades incluidas} \\
\midrule
\endhead
Periodos y categorías documentales (RF-01) &
Crear, listar, editar, cerrar y eliminar periodos; crear, editar y eliminar categorías
documentales; reutilizar categorías de periodos anteriores al crear un periodo nuevo.
Validaciones: nombre de periodo único, fecha de fin posterior a la de inicio, categoría
con nombre y no eliminar una categoría que ya tiene documentos. \\
\addlinespace
Facultades, carreras y usuarios (RF-02) &
Crear, listar, editar y eliminar facultades, carreras (cada una pertenece a una facultad)
y usuarios (con rol y carrera). Validaciones: nombre único, correo con formato válido y
único, carrera con facultad, carrera obligatoria salvo para el administrador general, y
no eliminar una facultad con carreras ni un usuario con documentos. \\
\addlinespace
Documentos (RF-03) &
Subir documentos en PDF o DOCX, consultarlos con filtros por periodo, categoría y carrera
con paginación y orden por fecha, editar sus datos, eliminarlos y descargarlos. \\
\bottomrule
\end{longtable}

Además, el alcance incluye la conexión de la aplicación con PostgreSQL, los scripts SQL
por requisito (\texttt{00\_base.sql}, \texttt{rf01.sql}, \texttt{rf02.sql} y
\texttt{rf03.sql}), las pruebas funcionales con evidencia y un script de datos de prueba
masivos para comprobar el rendimiento del diseño.

\subsection{Qué queda fuera}

\begin{itemize}[leftmargin=1.2cm]
  \item Inicio de sesión y control de acceso, que se contemplan para una versión posterior.
        Por ahora los roles se registran en el sistema, pero todavía no se restringen las
        pantallas según el rol de quien entra.
  \item Administración avanzada de roles y privilegios de PostgreSQL, que corresponde a
        contenidos posteriores de la materia.
  \item Carga real de millones de documentos. Solo se generan datos ficticios para
        probar que el diseño responde bien.
  \item Videollamadas, mensajería entre usuarios y agendamiento de tutorías, porque el
        problema que se resuelve es la dispersión y la pérdida de los documentos.
  \item Reportes y estadísticas de uso de la plataforma.
  \item Avisos o recordatorios por correo a los tutores.
  \item Importar automáticamente los documentos históricos que hoy están en correos.
\end{itemize}

\subsection{Cambios respecto a la propuesta inicial}

Después de la propuesta de la semana 2 el docente precisó cómo debe funcionar la
plataforma. Por eso se ajustó la propuesta inicial. La Tabla~\ref{tab:cambios} resume
los cambios y su motivo.

\begingroup\small
\begin{longtable}{>{\raggedright\arraybackslash}p{2.3cm}>{\raggedright\arraybackslash}p{3.6cm}>{\raggedright\arraybackslash}p{4.3cm}>{\raggedright\arraybackslash}p{3.9cm}}
\caption{Cambios respecto a la propuesta de la semana 2}\label{tab:cambios}\\
\toprule
\textbf{Elemento} & \textbf{Propuesta inicial} & \textbf{Versión actual} & \textbf{Motivo} \\
\midrule
\endfirsthead
\toprule
\textbf{Elemento} & \textbf{Propuesta inicial} & \textbf{Versión actual} & \textbf{Motivo} \\
\midrule
\endhead
Problema &
Materiales de tutorías para estudiantes y docentes. &
Documentos de tutorías (plan de acción inicial, informe de corte e informe final)
que hoy se buscan en años de correos. &
El docente explicó que el objetivo es evitar revisar cinco años de correos. \\
\addlinespace
Actores &
Administrador, docente y estudiante (tres actores). &
Administrador general, coordinador de carrera, coordinador de área y tutor (cuatro
roles). El estudiante ya no es un actor. &
El docente indicó estos roles; los tutores son quienes suben los documentos. \\
\addlinespace
Cobertura &
No se precisaba. &
Cualquier facultad y carrera, con facultades y carreras registradas como datos. &
El docente pidió que no sirva solo a la carrera de Software. \\
\addlinespace
Organización de los documentos &
Por asignatura, tema, tipo y fecha. &
Por periodo académico y categoría documental, reutilizable entre semestres. &
El administrador define cada semestre qué documentos se necesitan y las categorías
se reutilizan. \\
\addlinespace
Seguridad &
Inicio de sesión, control de acceso, cierre de sesión y contraseñas cifradas. &
Se deja para una versión posterior. Por ahora solo se registra el rol de cada usuario. &
Decisión del equipo para priorizar los tres módulos; la administración avanzada de
roles corresponde a contenidos posteriores. \\
\addlinespace
Rendimiento &
Búsquedas ágiles con al menos 1\,000 recursos. &
Diseño pensado para millones de registros, con índices, paginación y datos de prueba
masivos. &
El docente indicó que las bases de datos de la materia manejan millones de registros. \\
\bottomrule
\end{longtable}
\endgroup

% Paquetes que usa: booktabs, longtable, array, enumitem (ya van en el preambulo habitual).
% =====================================================================

\section{Misión}

El \textit{Repositorio de Tutorías} existe para guardar en un solo lugar, y clasificados
por periodo académico y por categoría documental, los documentos de las tutorías: el plan
de acción inicial, el informe de corte y el informe final.

Hoy esa información está repartida en correos electrónicos, archivos sueltos y
dispositivos personales. Cuando a un coordinador le piden un histórico de cinco años,
tiene que revisar cinco años de correos, y al no existir un lugar centralizado se pierde
información y cuesta encontrarla. Con la plataforma basta con cambiar de periodo y
consultar lo que se subió en ese semestre, en cualquier momento y desde cualquier lugar.

La plataforma no está pensada solo para la carrera de Software. Las facultades y las
carreras se registran como datos del sistema, de modo que cualquier facultad y cualquier
carrera de la institución pueda usarla con su propia información.

\section{Stakeholders}

La Tabla~\ref{tab:stakeholders} indica quiénes tienen interés en la plataforma o se ven
afectados por ella, y cuál es su relación con el sistema. Los tutores son quienes suben
los documentos; los demás roles organizan, consultan o evalúan.

\begin{longtable}{>{\raggedright\arraybackslash}p{3.7cm}>{\raggedright\arraybackslash}p{11.3cm}}
\caption{Stakeholders del sistema y su relación con él}\label{tab:stakeholders}\\
\toprule
\textbf{Stakeholder} & \textbf{Relación con el sistema} \\
\midrule
\endfirsthead
\toprule
\textbf{Stakeholder} & \textbf{Relación con el sistema} \\
\midrule
\endhead
Administrador general &
Configura la plataforma: registra facultades, carreras y usuarios, crea los periodos
y define qué categorías documentales se piden en cada semestre (por ejemplo, plan de
acción inicial, informe de corte e informe final). \\
\addlinespace
Coordinador de carrera &
Da seguimiento a los documentos de su carrera. Consulta lo que entregaron los tutores,
filtrando por periodo y categoría. \\
\addlinespace
Coordinador de área &
Revisa los documentos de las carreras de su área. \\
\addlinespace
Tutores &
Son quienes suben los documentos. Eligen el periodo y la categoría, escriben el título
y adjuntan el archivo, que queda guardado en PostgreSQL. \\
\addlinespace
Docente de la materia (evaluador) &
Orienta y evalúa el proyecto. Revisa que la aplicación funcione, que la integración con
PostgreSQL sea correcta y que las decisiones estén documentadas. \\
\addlinespace
Institución (UTEQ) &
Es la beneficiaria final: conserva el histórico de tutorías ordenado y accesible,
sin depender de correos personales. \\
\bottomrule
\end{longtable}

\section{Objetivos}

\subsection{Objetivo general}

Desarrollar una aplicación web conectada a PostgreSQL que guarde y clasifique por
periodo académico y categoría documental los documentos de tutorías, para que pueda
ser usada por cualquier facultad y carrera y reemplace la búsqueda de documentos en
años de correos.

\subsection{Objetivos específicos}

\begin{enumerate}[label=\textbf{OE\arabic*.}, leftmargin=1.6cm]
  \item Gestionar periodos y categorías documentales, de modo que al crear un periodo
        nuevo se puedan reutilizar las categorías de periodos anteriores y agregar
        nuevas o quitar las que ya no se necesiten (módulo de periodos y categorías,
        RF-01).
  \item Registrar facultades, carreras y usuarios, asignando a cada usuario uno de los
        cuatro roles (administrador general, coordinador de carrera, coordinador de área
        o tutor), con validaciones de integridad como nombre único, correo válido y
        único, y carrera obligatoria para todos los roles excepto el administrador
        general (módulo de facultades, carreras y usuarios, RF-02).
  \item Permitir que los tutores suban documentos y que se consulten con filtros por
        periodo, categoría y carrera, con paginación, validando el tipo de archivo, el
        tamaño y que la categoría esté habilitada para ese periodo (módulo de
        documentos, RF-03).
  \item Diseñar la base de datos para que soporte millones de registros, usando índices
        en las llaves foráneas y en los campos de filtro, paginación en las listas y un
        script de datos de prueba masivos, para que las pantallas no se vuelvan lentas
        al crecer la información.
  \item Proteger la integridad de los datos con llaves foráneas \texttt{ON DELETE
        RESTRICT}, consultas parametrizadas y mensajes claros de éxito o de error en
        cada operación.
  \item Comprobar cada requisito con pruebas funcionales que registren los datos
        usados, el resultado esperado, el resultado obtenido y su evidencia, y que
        verifiquen en PostgreSQL lo que se hizo desde la aplicación.
\end{enumerate}

\section{Alcance}

\subsection{Qué incluye este semestre}

La primera versión de la aplicación es una aplicación web en C\# con ASP.NET Core MVC
y PostgreSQL como gestor de base de datos, donde se almacena toda la información. Incluye
tres módulos de prioridad alta (Tabla~\ref{tab:modulos}):

\begin{longtable}{>{\raggedright\arraybackslash}p{3.7cm}>{\raggedright\arraybackslash}p{11.3cm}}
\caption{Módulos incluidos en la primera versión}\label{tab:modulos}\\
\toprule
\textbf{Módulo} & \textbf{Funcionalidades incluidas} \\
\midrule
\endfirsthead
\toprule
\textbf{Módulo} & \textbf{Funcionalidades incluidas} \\
\midrule
\endhead
Periodos y categorías documentales (RF-01) &
Crear, listar, editar, cerrar y eliminar periodos; crear, editar y eliminar categorías
documentales; reutilizar categorías de periodos anteriores al crear un periodo nuevo.
Validaciones: nombre de periodo único, fecha de fin posterior a la de inicio, categoría
con nombre y no eliminar una categoría que ya tiene documentos. \\
\addlinespace
Facultades, carreras y usuarios (RF-02) &
Crear, listar, editar y eliminar facultades, carreras (cada una pertenece a una facultad)
y usuarios (con rol y carrera). Validaciones: nombre único, correo con formato válido y
único, carrera con facultad, carrera obligatoria salvo para el administrador general, y
no eliminar una facultad con carreras ni un usuario con documentos. \\
\addlinespace
Documentos (RF-03) &
Subir documentos en PDF o DOCX, consultarlos con filtros por periodo, categoría y carrera
con paginación y orden por fecha, editar sus datos, eliminarlos y descargarlos. \\
\bottomrule
\end{longtable}

Además, el alcance incluye la conexión de la aplicación con PostgreSQL, los scripts SQL
por requisito (\texttt{00\_base.sql}, \texttt{rf01.sql}, \texttt{rf02.sql} y
\texttt{rf03.sql}), las pruebas funcionales con evidencia y un script de datos de prueba
masivos para comprobar el rendimiento del diseño.

\subsection{Qué queda fuera}

\begin{itemize}[leftmargin=1.2cm]
  \item Inicio de sesión y control de acceso, que se contemplan para una versión posterior.
        Por ahora los roles se registran en el sistema, pero todavía no se restringen las
        pantallas según el rol de quien entra.
  \item Administración avanzada de roles y privilegios de PostgreSQL, que corresponde a
        contenidos posteriores de la materia.
  \item Carga real de millones de documentos. Solo se generan datos ficticios para
        probar que el diseño responde bien.
  \item Videollamadas, mensajería entre usuarios y agendamiento de tutorías, porque el
        problema que se resuelve es la dispersión y la pérdida de los documentos.
  \item Reportes y estadísticas de uso de la plataforma.
  \item Avisos o recordatorios por correo a los tutores.
  \item Importar automáticamente los documentos históricos que hoy están en correos.
\end{itemize}

\subsection{Cambios respecto a la propuesta inicial}

Después de la propuesta de la semana 2 el docente precisó cómo debe funcionar la
plataforma. Por eso se ajustó la propuesta inicial. La Tabla~\ref{tab:cambios} resume
los cambios y su motivo.

\begingroup\small
\begin{longtable}{>{\raggedright\arraybackslash}p{2.3cm}>{\raggedright\arraybackslash}p{3.6cm}>{\raggedright\arraybackslash}p{4.3cm}>{\raggedright\arraybackslash}p{3.9cm}}
\caption{Cambios respecto a la propuesta de la semana 2}\label{tab:cambios}\\
\toprule
\textbf{Elemento} & \textbf{Propuesta inicial} & \textbf{Versión actual} & \textbf{Motivo} \\
\midrule
\endfirsthead
\toprule
\textbf{Elemento} & \textbf{Propuesta inicial} & \textbf{Versión actual} & \textbf{Motivo} \\
\midrule
\endhead
Problema &
Materiales de tutorías para estudiantes y docentes. &
Documentos de tutorías (plan de acción inicial, informe de corte e informe final)
que hoy se buscan en años de correos. &
El docente explicó que el objetivo es evitar revisar cinco años de correos. \\
\addlinespace
Actores &
Administrador, docente y estudiante (tres actores). &
Administrador general, coordinador de carrera, coordinador de área y tutor (cuatro
roles). El estudiante ya no es un actor. &
El docente indicó estos roles; los tutores son quienes suben los documentos. \\
\addlinespace
Cobertura &
No se precisaba. &
Cualquier facultad y carrera, con facultades y carreras registradas como datos. &
El docente pidió que no sirva solo a la carrera de Software. \\
\addlinespace
Organización de los documentos &
Por asignatura, tema, tipo y fecha. &
Por periodo académico y categoría documental, reutilizable entre semestres. &
El administrador define cada semestre qué documentos se necesitan y las categorías
se reutilizan. \\
\addlinespace
Seguridad &
Inicio de sesión, control de acceso, cierre de sesión y contraseñas cifradas. &
Se deja para una versión posterior. Por ahora solo se registra el rol de cada usuario. &
Decisión del equipo para priorizar los tres módulos; la administración avanzada de
roles corresponde a contenidos posteriores. \\
\addlinespace
Rendimiento &
Búsquedas ágiles con al menos 1\,000 recursos. &
Diseño pensado para millones de registros, con índices, paginación y datos de prueba
masivos. &
El docente indicó que las bases de datos de la materia manejan millones de registros. \\
\bottomrule
\end{longtable}
\endgroup

% Paquetes que usa: booktabs, longtable, array, enumitem (ya van en el preambulo habitual).
% =====================================================================

\section{Misión}

El \textit{Repositorio de Tutorías} existe para guardar en un solo lugar, y clasificados
por periodo académico y por categoría documental, los documentos de las tutorías: el plan
de acción inicial, el informe de corte y el informe final.

Hoy esa información está repartida en correos electrónicos, archivos sueltos y
dispositivos personales. Cuando a un coordinador le piden un histórico de cinco años,
tiene que revisar cinco años de correos, y al no existir un lugar centralizado se pierde
información y cuesta encontrarla. Con la plataforma basta con cambiar de periodo y
consultar lo que se subió en ese semestre, en cualquier momento y desde cualquier lugar.

La plataforma no está pensada solo para la carrera de Software. Las facultades y las
carreras se registran como datos del sistema, de modo que cualquier facultad y cualquier
carrera de la institución pueda usarla con su propia información.

\section{Stakeholders}

La Tabla~\ref{tab:stakeholders} indica quiénes tienen interés en la plataforma o se ven
afectados por ella, y cuál es su relación con el sistema. Los tutores son quienes suben
los documentos; los demás roles organizan, consultan o evalúan.

\begin{longtable}{>{\raggedright\arraybackslash}p{3.7cm}>{\raggedright\arraybackslash}p{11.3cm}}
\caption{Stakeholders del sistema y su relación con él}\label{tab:stakeholders}\\
\toprule
\textbf{Stakeholder} & \textbf{Relación con el sistema} \\
\midrule
\endfirsthead
\toprule
\textbf{Stakeholder} & \textbf{Relación con el sistema} \\
\midrule
\endhead
Administrador general &
Configura la plataforma: registra facultades, carreras y usuarios, crea los periodos
y define qué categorías documentales se piden en cada semestre (por ejemplo, plan de
acción inicial, informe de corte e informe final). \\
\addlinespace
Coordinador de carrera &
Da seguimiento a los documentos de su carrera. Consulta lo que entregaron los tutores,
filtrando por periodo y categoría. \\
\addlinespace
Coordinador de área &
Revisa los documentos de las carreras de su área. \\
\addlinespace
Tutores &
Son quienes suben los documentos. Eligen el periodo y la categoría, escriben el título
y adjuntan el archivo, que queda guardado en PostgreSQL. \\
\addlinespace
Docente de la materia (evaluador) &
Orienta y evalúa el proyecto. Revisa que la aplicación funcione, que la integración con
PostgreSQL sea correcta y que las decisiones estén documentadas. \\
\addlinespace
Institución (UTEQ) &
Es la beneficiaria final: conserva el histórico de tutorías ordenado y accesible,
sin depender de correos personales. \\
\bottomrule
\end{longtable}

\section{Objetivos}

\subsection{Objetivo general}

Desarrollar una aplicación web conectada a PostgreSQL que guarde y clasifique por
periodo académico y categoría documental los documentos de tutorías, para que pueda
ser usada por cualquier facultad y carrera y reemplace la búsqueda de documentos en
años de correos.

\subsection{Objetivos específicos}

\begin{enumerate}[label=\textbf{OE\arabic*.}, leftmargin=1.6cm]
  \item Gestionar periodos y categorías documentales, de modo que al crear un periodo
        nuevo se puedan reutilizar las categorías de periodos anteriores y agregar
        nuevas o quitar las que ya no se necesiten (módulo de periodos y categorías,
        RF-01).
  \item Registrar facultades, carreras y usuarios, asignando a cada usuario uno de los
        cuatro roles (administrador general, coordinador de carrera, coordinador de área
        o tutor), con validaciones de integridad como nombre único, correo válido y
        único, y carrera obligatoria para todos los roles excepto el administrador
        general (módulo de facultades, carreras y usuarios, RF-02).
  \item Permitir que los tutores suban documentos y que se consulten con filtros por
        periodo, categoría y carrera, con paginación, validando el tipo de archivo, el
        tamaño y que la categoría esté habilitada para ese periodo (módulo de
        documentos, RF-03).
  \item Diseñar la base de datos para que soporte millones de registros, usando índices
        en las llaves foráneas y en los campos de filtro, paginación en las listas y un
        script de datos de prueba masivos, para que las pantallas no se vuelvan lentas
        al crecer la información.
  \item Proteger la integridad de los datos con llaves foráneas \texttt{ON DELETE
        RESTRICT}, consultas parametrizadas y mensajes claros de éxito o de error en
        cada operación.
  \item Comprobar cada requisito con pruebas funcionales que registren los datos
        usados, el resultado esperado, el resultado obtenido y su evidencia, y que
        verifiquen en PostgreSQL lo que se hizo desde la aplicación.
\end{enumerate}

\section{Alcance}

\subsection{Qué incluye este semestre}

La primera versión de la aplicación es una aplicación web en C\# con ASP.NET Core MVC
y PostgreSQL como gestor de base de datos, donde se almacena toda la información. Incluye
tres módulos de prioridad alta (Tabla~\ref{tab:modulos}):

\begin{longtable}{>{\raggedright\arraybackslash}p{3.7cm}>{\raggedright\arraybackslash}p{11.3cm}}
\caption{Módulos incluidos en la primera versión}\label{tab:modulos}\\
\toprule
\textbf{Módulo} & \textbf{Funcionalidades incluidas} \\
\midrule
\endfirsthead
\toprule
\textbf{Módulo} & \textbf{Funcionalidades incluidas} \\
\midrule
\endhead
Periodos y categorías documentales (RF-01) &
Crear, listar, editar, cerrar y eliminar periodos; crear, editar y eliminar categorías
documentales; reutilizar categorías de periodos anteriores al crear un periodo nuevo.
Validaciones: nombre de periodo único, fecha de fin posterior a la de inicio, categoría
con nombre y no eliminar una categoría que ya tiene documentos. \\
\addlinespace
Facultades, carreras y usuarios (RF-02) &
Crear, listar, editar y eliminar facultades, carreras (cada una pertenece a una facultad)
y usuarios (con rol y carrera). Validaciones: nombre único, correo con formato válido y
único, carrera con facultad, carrera obligatoria salvo para el administrador general, y
no eliminar una facultad con carreras ni un usuario con documentos. \\
\addlinespace
Documentos (RF-03) &
Subir documentos en PDF o DOCX, consultarlos con filtros por periodo, categoría y carrera
con paginación y orden por fecha, editar sus datos, eliminarlos y descargarlos. \\
\bottomrule
\end{longtable}

Además, el alcance incluye la conexión de la aplicación con PostgreSQL, los scripts SQL
por requisito (\texttt{00\_base.sql}, \texttt{rf01.sql}, \texttt{rf02.sql} y
\texttt{rf03.sql}), las pruebas funcionales con evidencia y un script de datos de prueba
masivos para comprobar el rendimiento del diseño.

\subsection{Qué queda fuera}

\begin{itemize}[leftmargin=1.2cm]
  \item Inicio de sesión y control de acceso, que se contemplan para una versión posterior.
        Por ahora los roles se registran en el sistema, pero todavía no se restringen las
        pantallas según el rol de quien entra.
  \item Administración avanzada de roles y privilegios de PostgreSQL, que corresponde a
        contenidos posteriores de la materia.
  \item Carga real de millones de documentos. Solo se generan datos ficticios para
        probar que el diseño responde bien.
  \item Videollamadas, mensajería entre usuarios y agendamiento de tutorías, porque el
        problema que se resuelve es la dispersión y la pérdida de los documentos.
  \item Reportes y estadísticas de uso de la plataforma.
  \item Avisos o recordatorios por correo a los tutores.
  \item Importar automáticamente los documentos históricos que hoy están en correos.
\end{itemize}

\subsection{Cambios respecto a la propuesta inicial}

Después de la propuesta de la semana 2 el docente precisó cómo debe funcionar la
plataforma. Por eso se ajustó la propuesta inicial. La Tabla~\ref{tab:cambios} resume
los cambios y su motivo.

\begingroup\small
\begin{longtable}{>{\raggedright\arraybackslash}p{2.3cm}>{\raggedright\arraybackslash}p{3.6cm}>{\raggedright\arraybackslash}p{4.3cm}>{\raggedright\arraybackslash}p{3.9cm}}
\caption{Cambios respecto a la propuesta de la semana 2}\label{tab:cambios}\\
\toprule
\textbf{Elemento} & \textbf{Propuesta inicial} & \textbf{Versión actual} & \textbf{Motivo} \\
\midrule
\endfirsthead
\toprule
\textbf{Elemento} & \textbf{Propuesta inicial} & \textbf{Versión actual} & \textbf{Motivo} \\
\midrule
\endhead
Problema &
Materiales de tutorías para estudiantes y docentes. &
Documentos de tutorías (plan de acción inicial, informe de corte e informe final)
que hoy se buscan en años de correos. &
El docente explicó que el objetivo es evitar revisar cinco años de correos. \\
\addlinespace
Actores &
Administrador, docente y estudiante (tres actores). &
Administrador general, coordinador de carrera, coordinador de área y tutor (cuatro
roles). El estudiante ya no es un actor. &
El docente indicó estos roles; los tutores son quienes suben los documentos. \\
\addlinespace
Cobertura &
No se precisaba. &
Cualquier facultad y carrera, con facultades y carreras registradas como datos. &
El docente pidió que no sirva solo a la carrera de Software. \\
\addlinespace
Organización de los documentos &
Por asignatura, tema, tipo y fecha. &
Por periodo académico y categoría documental, reutilizable entre semestres. &
El administrador define cada semestre qué documentos se necesitan y las categorías
se reutilizan. \\
\addlinespace
Seguridad &
Inicio de sesión, control de acceso, cierre de sesión y contraseñas cifradas. &
Se deja para una versión posterior. Por ahora solo se registra el rol de cada usuario. &
Decisión del equipo para priorizar los tres módulos; la administración avanzada de
roles corresponde a contenidos posteriores. \\
\addlinespace
Rendimiento &
Búsquedas ágiles con al menos 1\,000 recursos. &
Diseño pensado para millones de registros, con índices, paginación y datos de prueba
masivos. &
El docente indicó que las bases de datos de la materia manejan millones de registros. \\
\bottomrule
\end{longtable}
\endgroup
